\section*{Chapitre \ref{ch:b3:relativistic-kinematics} --- Cinématique relativiste}

\begin{solution}{exo:b3:relativistic-kinematics:1}
$\gamma = 1.005$, $1.155$, $2.294$, $7.09$, $22.4$. Station spatiale~: $\beta =
\num{2.57e-5}$, $\gamma - 1 \approx \beta^2/2 = \num{3.3e-10}$~; sur
six mois ($\num{1.6e7}{}\,\unit{s}$) l'horloge de l'équipage
\emph{retarde} de $\approx\qty{5}{ms}$ (effet de vitesse seul~; à la
faible altitude de la station, l'effet gravitationnel le réduit sans
l'inverser).
\end{solution}

\begin{solution}{exo:b3:relativistic-kinematics:2}
(a) $\Delta\tau = 2d/c$. (b) À chaque demi-battement, le photon
parcourt l'hypoténuse~: $(c\Delta t/2)^2 = (v\Delta t/2)^2 + d^2$ avec $d =
c\Delta\tau/2$~; on résout~: $\Delta t = \Delta\tau/\sqrt{1 - v^2/c^2}$. (c)
Si $d$ changeait avec le mouvement, l'étape de Pythagore serait fausse.
(d) Faisons se croiser deux règles identiques, chacune portant à son
extrémité un pinceau dirigé vers l'autre. Si le mouvement contractait
les longueurs transverses, chaque référentiel prédirait que sa propre
règle marque \emph{au-delà} de l'extrémité de l'autre --- deux faits
contradictoires sur les mêmes traits de pinceau, au même \omterm{def:b3:relativistic-kinematics:event}{événement} de
croisement. Une contradiction en un même \omterm{def:b3:relativistic-kinematics:event}{événement} est interdite~: les
longueurs transverses ne peuvent pas changer.
\end{solution}

\begin{solution}{exo:b3:relativistic-kinematics:3}
$\gamma = 22.4$. (a) $\gamma\tau = \qty{49}{\micro s}$. (b) $v\gamma
\tau = \qty{14.8}{km}$. (c) $15/22.4 = \qty{670}{m}$. (d) Temps de vol
au laboratoire $\qty{50}{\micro s}$~: sans la relativité
$\eu^{-50/2.2} \approx \num{e-10}$ --- aucun~; avec la relativité le
\omterm{prop:b3:relativistic-kinematics:dilation}{temps propre} vaut $50/22.4 = \qty{2.2}{\micro s}$, fraction $\eu^{-1}
\approx 0.37$.
\end{solution}

\begin{solution}{exo:b3:relativistic-kinematics:4}
(a) $c\Delta t = \qty{600}{m}$, $\Delta x = \qty{300}{m}$~: $\Delta s^2
= (600^2 - 300^2)\,\unit{m^2} > 0$, genre temps. (b) Oui~: un signal à
$\Delta x/\Delta t = c/2$ suffit. (c) Ce même référentiel, $v =
0.5c$~; \omterm{prop:b3:relativistic-kinematics:dilation}{temps propre} $\Delta s/c = \sqrt{600^2 - 300^2}/c = 520/c =
\qty{1.73}{\micro s}$. (d) Maintenant $c\Delta t = 300 < \Delta x = 600$~:
genre espace~; aucun lien causal possible~; aucun référentiel ne les
amène au même lieu, mais celui de vitesse $v = c^2\Delta t/\Delta x = 0.5c$ les rend
simultanés.
\end{solution}

\begin{solution}{exo:b3:relativistic-kinematics:5}
(a) $\beta = \num{1.29e-5}$~: $\gamma - 1 = \num{8.3e-11}$. (b)
$86400\,\unit{s} \times \num{8.3e-11} = \qty{7.2}{\micro s}$ par jour.
(c) $c \times \qty{7.2}{\micro s} = \qty{2.2}{km}$ --- la navigation
mourrait en un jour. (d) Le solde $+38\,\unit{\micro s}$ signifie que
le décalage gravitationnel vers le bleu ($+45.7$) l'emporte sur le
ralentissement cinématique ($-7.2$)~: à \qty{20200}{km}, être plus haut
dans le potentiel terrestre accélère l'horloge davantage que la vitesse
orbitale ne la ralentit.
\end{solution}

\begin{solution}{exo:b3:relativistic-kinematics:6}
(a) $1.6c/1.64 = 0.976c$. (b) $1.8c/1.81 = 0.994c$. (c) $c - u =
\dfrac{(c - u')(c - v)}{c\,(1 + u'v/c^2)}$~: les deux facteurs sont
positifs, donc $u < c$. (d) La tache est un \emph{motif} en mouvement,
non une chose~: ni matière, ni énergie, ni information ne va d'un point
de la face lunaire au suivant --- chaque photon est parti vers la Lune
à $c$.
\end{solution}

\begin{solution}{exo:b3:relativistic-kinematics:7}
(a) Longueur d'onde plus courte~: elle s'approche. $f/f_0 = 589/575 = 1.024$~;
$(1+\beta)/(1-\beta) = 1.049$~: $\beta = 0.024$, soit environ
\qty{7200}{km/s}. (b) Rapport $2$~: $(1+\beta)/(1-\beta) = 4$, $\beta =
0.6$. (c) En développant la racine, $\Delta\lambda/\lambda \approx
\beta$~; à \qty{600}{nm}, $\qty{1}{\angstrom} = \qty{0.1}{nm}$ donne
$\beta = \num{1.7e-4}$~: environ \qty{50}{km/s} par ångström --- le
réflexe de l'astronome. (d) L'effet Doppler transverse~: pure
\omterm{prop:b3:relativistic-kinematics:dilation}{dilatation du temps}, d'ordre $\beta^2$ --- mesurable seulement avec une
précision atomique.
\end{solution}

\begin{solution}{exo:b3:relativistic-kinematics:8}
(a) $20/2 = \qty{10}{m}$~: elle y tient, tout juste, et les deux portes
peuvent être fermées simultanément (référentiel de la grange). (b) Dans
le référentiel du coureur, les \omterm{def:b3:relativistic-kinematics:event}{événements} «~la porte avant se ferme~»
et «~la porte arrière s'ouvre~» ne sont \emph{pas} simultanés~: la
sortie s'ouvre avant que l'entrée ne se ferme, et la perche de
\qty{20}{m} enfile une grange de \qty{5}{m} sans y être jamais
enfermée. (c) Référentiel de la grange~: $\Delta t = 0$ sur $\Delta x =
\qty{10}{m}$. Référentiel du coureur~: $|\Delta t'| = \gamma v\Delta x/c^2 =
2 \times 0.866c \times 10/c^2 = \qty{58}{ns}$. (d) «~La perche est
enfermée~» affirme tacitement que deux \omterm{def:b3:relativistic-kinematics:event}{événements} séparés (les deux
portes fermées) sont simultanés --- un énoncé lié au référentiel, non
un fait.
\end{solution}

\begin{solution}{exo:b3:relativistic-kinematics:9}
$\gamma = 1.25$ à $0.6c$. (a) Longueur en mouvement $100/1.25 = \qty{80}{m}$
à \qty{1.8e8}{m/s}~: \qty{444}{ns}. (b) $240/\num{1.8e8} =
\qty{1.33}{\micro s}$. (c) $|\Delta t'| = \gamma v\Delta x/c^2 =
\qty{200}{ns}$~; d'après $t' = \gamma(t - vx/c^2)$, le verrou de plus
grand $x$ --- celui de l'avant --- part le \emph{premier} pour la
fusée. (d) Station~: $\Delta s^2 = 0 - 80^2 = -6400\,\unit{m^2}$. Fusée~: $\Delta
x' = \gamma\,\Delta x = \qty{100}{m}$, $c\Delta t' = \qty{60}{m}$~:
$60^2 - 100^2 = -6400\,\unit{m^2}$ --- invariant.
\end{solution}

\begin{solution}{exo:b3:relativistic-kinematics:10}
(a) Terra~: \qty{20}{yr}~; Stella~: $20/\gamma = \qty{12}{yr}$. (b) La
distance se contracte à $8/\gamma = \qty{4.8}{ly}$, parcourus en $4.8/0.8
= \qty{6}{yr}$ de son temps --- cohérent avec les \qty{12}{yr} de
l'aller-retour. (c) \omterm{def:b3:relativistic-kinematics:event}{Simultanéité} avec l'\omterm{def:b3:relativistic-kinematics:event}{événement} du demi-tour $(t =
\qty{10}{yr},\ x = \qty{8}{ly})$~: le référentiel d'aller attribue à
l'horloge terrestre $t - vx/c^2 = 10 - 6.4 = \qty{3.6}{yr}$~; celui du retour, $10 +
6.4 = \qty{16.4}{yr}$. Son «~maintenant sur Terre~» saute de \qty{12.8}{yr}
au coude~; son bilan $3.6 + 12.8 + 3.6 = \qty{20}{yr}$
retrouve exactement Terra. (d) Stella occupe \emph{deux} \omterm{thm:b3:relativistic-kinematics:postulates}{référentiels
galiléens} raccordés par une accélération qu'elle ressent~; Terra n'en
occupe qu'un. La ligne d'univers droite entre les \omterm{def:b3:relativistic-kinematics:event}{événements} de départ
et de retrouvailles porte le plus long \omterm{prop:b3:relativistic-kinematics:dilation}{temps propre}~; seule celle de
Stella est brisée.
\end{solution}

\begin{solution}{exo:b3:relativistic-kinematics:11}
(a) Si l'application était non linéaire, un mouvement uniforme ne
paraîtrait pas uniforme dans l'autre référentiel, ce qui distinguerait
des points d'un espace et d'un temps homogènes. (b) En reportant une
relation dans l'autre~: $t' = \gamma t +
(1 - \gamma^2)x/\gamma v$. (c) En posant $x = ct$, $x' = ct'$~:
$\gamma(c - v)t = c\gamma t + c(1 - \gamma^2)ct/\gamma v$, ce qui donne
$\gamma^2 = 1/(1 - v^2/c^2)$. (d) La relativité a donné le même
$\gamma$ pour les deux sens (aucun référentiel privilégié)~; le
postulat de la lumière a changé la fonction libre restante en le
$\gamma(v)$ déterminé.
\end{solution}

\begin{solution}{exo:b3:relativistic-kinematics:12}
(a) Avec $\cosh\varphi = \gamma$, $\sinh\varphi = \gamma\beta$, la
transformation est exactement la rotation hyperbolique annoncée, et
$\cosh^2 - \sinh^2 = 1$ conserve $c^2t^2 - x^2$. (b)
$\tanh(\varphi_1 + \varphi_2) = (\tanh\varphi_1 + \tanh\varphi_2)/(1 +
\tanh\varphi_1\tanh\varphi_2)$~: précisément la loi de composition ---
les vitesses ne s'ajoutent pas, les rapidités si. (c) $\varphi = g\tau/c$, donc
$\beta = \tanh(g\tau/c)$~; $\operatorname{artanh}(0.99) = 2.65$ donne
$\tau = 2.65c/g \approx \qty{8.1e7}{s} \approx 2.6$ ans de temps de
bord. (d) $\varphi$ parcourt tous les réels tandis que $\tanh\varphi$
sature à $1$~: on peut accélérer indéfiniment, gagnant de la rapidité
linéairement, pendant que la vitesse ne fait que ramper vers $c$.
\end{solution}

\begin{solution}{pb:b3:relativistic-kinematics:1}
\textbf{1.} Tous les muons sont rigoureusement identiques et se
désintègrent à un taux statistique fixe~: la fraction survivante d'une
population mesure le \omterm{prop:b3:relativistic-kinematics:dilation}{temps propre} écoulé aussi sûrement qu'une aiguille.
\textbf{2.} $c\tau = \qty{3.00e8}{} \times \qty{2.2e-6}{} =
\qty{660}{m}$.
\textbf{3.} $t = \qty{15}{km}/c = \qty{50}{\micro s}$~: fraction
$\eu^{-50/2.2} = \eu^{-22.7} \approx \num{1.4e-10}$.
\textbf{4.} Des particules qui «~ne peuvent pas~» franchir plus d'un
kilomètre en traversent quinze et arrivent en nombre.
\textbf{5.} Une main de $\sim\qty{100}{cm^2}$~: deux ou trois muons par
seconde~; un corps de section horizontale $\sim\num{e3}\,\unit{cm^2}$~:
une quinzaine entre deux battements de cœur --- la relativité pleut sur
tout le monde, en permanence.
\textbf{6.} $\gamma = 1/\sqrt{1 - 0.995^2} = 10.0$.
\textbf{7.} Comparer deux comptages de la \emph{même} population
sélectionnée sur un dénivelé connu élimine toute hypothèse sur le lieu
et le nombre des naissances de muons.
\textbf{8.} $t = 1907/(0.995 \times \num{3e8}) = \qty{6.39}{\micro s}$.
\textbf{9.} $\eu^{-6.39/2.2} = 0.055$~: environ $31$ par heure.
\textbf{10.} $t/\gamma = \qty{0.64}{\micro s}$.
\textbf{11.} $\eu^{-0.64/2.2} = 0.75$~: environ $420$ par heure.
\textbf{12.} Mesure $408 \pm 9$~: les $\approx 420$ de la relativité
s'accordent, au léger ralentissement des muons dans l'absorbeur des
détecteurs près~; les $31$ de la physique classique sont faux d'un
facteur treize. La montagne tranche.
\textbf{13.} $408/563 = 0.725 = \eu^{-t_{\text{propre}}/\tau}$~:
$t_{\text{propre}} = \qty{0.71}{\micro s}$, donc $\gamma_{\exp} =
6.39/0.71 = 9.0$ --- en plein dans le $8.8 \pm 0.8$ de Frisch et Smith.
\textbf{14.} $1907/10 = \qty{191}{m}$.
\textbf{15.} $191/(0.995c) = \qty{0.64}{\micro s}$~: les mêmes $25\%$
de désintégration --- la version du muon d'un même comptage.
\textbf{16.} Un comptage de clics est un ensemble d'\omterm{def:b3:relativistic-kinematics:event}{événements} locaux
de coïncidence~; les \omterm{def:b3:relativistic-kinematics:event}{événements} et leurs décomptes sont invariants ---
les référentiels peuvent différer sur les temps et les longueurs,
jamais sur ce qu'un compteur a affiché.
\textbf{17.} $\Delta s^2 = (c \times \qty{6.39}{\micro s})^2 -
(\qty{1907}{m})^2 = (1917^2 - 1907^2)\,\unit{m^2} = (\qty{196}{m})^2$~:
$\Delta s/c = \qty{0.65}{\micro s}$ --- le \omterm{prop:b3:relativistic-kinematics:dilation}{temps propre}, calculé sans
jamais quitter le référentiel de la montagne.
\textbf{18.} Les taux de désintégration dépendraient de l'altitude pour
toutes les vitesses de la même façon~; au lieu de cela la survie suit
$\gamma$ --- les sélections plus lentes se dilatent moins, exactement
comme $\gamma(v)$ le prescrit.
\textbf{19.} Des durées de vie de $\gamma\tau$ au millième près
(anneaux de stockage de muons du CERN)~; l'anneau fait du muon un
voyageur perpétuellement accéléré --- la «~jumelle voyageuse~» reste
jeune même quand le voyage n'est qu'un demi-tour sans fin.
\textbf{20.} $\beta = \num{8.3e-7}$~: $\Delta t = \tfrac12\beta^2
\times \qty{144000}{s} = \qty{50}{ns}$.
\textbf{21.} La vitesse de l'avion s'ajoute à celle de la Terre en
rotation ou s'en retranche, tandis que l'effet d'altitude
(gravitationnel) est le même dans les deux sens~: les vols vers l'est
et vers l'ouest séparent proprement les deux contributions (le
résultat de 1971 s'accordait aux deux).
\textbf{22.} D'après \cref{exo:b3:relativistic-kinematics:5}~:
\qty{7.2}{\micro s} par jour.
\textbf{23.} $7 \times 7.2\,\unit{\micro s} \times c \approx
\qty{15}{km}.$
\textbf{24.} Les positions dériveraient de centaines de mètres en
quelques heures, de kilomètres en un jour --- la navigation se
casserait visiblement dès le premier matin.
\textbf{25.} En 1963, une montagne, deux compteurs et des horloges de
\qty{2.2}{\micro s} ont mesuré $\gamma_{\exp} = 9.0$ contre un $10$
prédit (avec le ralentissement, $8.8 \pm 0.8$)~: la \omterm{prop:b3:relativistic-kinematics:dilation}{dilatation du temps}
confirmée à $10\%$ près --- et la physique identique est corrigée en
silence, chaque seconde, dans chaque satellite de navigation
au-dessus de votre tête.
\end{solution}
